\section*{Capítulo \ref{ch:g4:mult} --- Multiplicando números maiores}

\begin{solution}{exo:g4:mult:1}
$470$; \quad $2\,300$; \quad $720$; \quad $4\,500$; \quad
$2\,100$.
\end{solution}

\begin{solution}{exo:g4:mult:2}
$34 \times 21 = 34 \times 20 + 34 = 680 + 34 = 714$; as colunas
dê as mesmas linhas: $34$ e $680$.
\end{solution}

\begin{solution}{exo:g4:mult:3}
$63 \times 24 = 1\,512$; \quad
$85 \times 47 = 3\,995$; \quad
$126 \times 53 = 6\,678$.
\end{solution}

\begin{solution}{exo:g4:mult:4}
$208 \times 34$: estimativa $200 \times 34 = 6\,800$; exatamente
$7\,072$.

$517 \times 62$: estimativa $500 \times 60 = 30\,000$; exatamente
$32\,054$.
\end{solution}

\begin{solution}{exo:g4:mult:5}
$45 \times 11 = 450 + 45 = 495$.

$32 \times 99 = 3\,200 - 32 = 3\,168$.

$24 \times 45 = 12 \times 90 = 1\,080$ (ou mais uma vez:
$6 \times 180 = 1\,080$).
\end{solution}

\begin{solution}{exo:g4:mult:6}
$26 \times 18 = 468$ assentos.
\end{solution}

\begin{solution}{exo:g4:mult:7}
Estimativa: $30 \times 50 = 1\,500$. Exatamente:
$32 \times 48 = 1\,536$ caixas.
\end{solution}

\begin{solution}{exo:g4:mult:8}
Ela esqueceu que o $2$ de $23$ conta \emph{dezenas}: a segunda linha
deve ser $57 \times 20 = 1\,140$, não $114$. Total correto:
$171 + 1\,140 = 1\,311$.
\end{solution}

\begin{solution}{exo:g4:mult:9}
Tulipas: $14 \times 35 = 490$. Narcisos: $12 \times 28 = 336$.
Total: $490 + 336 = 826$ flores.
\end{solution}

\begin{solution}{exo:g4:mult:10}
Estimativa: $70 \times 30 = 2\,100$ kg. Exatamente:
$67 \times 30 = 2\,010$ kg.
\end{solution}

\begin{solution}{exo:g4:mult:11}
UM $15 \times 12$ retângulo de quadrados unitários pode ser cortado por uma vertical
linha em um $15 \times 10$ parte e um $15 \times 2$ parte: contando
cada parte dá $150 + 30 = 180$. Outro corte funciona tão bem,
por exemplo $15 = 10 + 5$: $10 \times 12 + 5 \times 12 = 120 + 60
= 180$. Qualquer forma de cortar o retângulo deixa o número total de
quadrados inalterados.
\end{solution}
